
% ============================================================================
\clearpage
\section*{v4.10.0-r54 Source-Only: Does the Readout Fix Its Own Frame? (H122, Q1b)}
\addcontentsline{toc}{section}{v4.10.0-r54 Source-Only: Does the Readout Fix Its Own Frame? (H122)}
\label{sec:v4100-r54-readout-frame}
% ============================================================================

\paragraph{Purpose.}
The cross-section Q1 of the open problems asks for one readout law that supplies the frame \(u\), the axis \(n\) and the time orientation used by
OP1 (reading B), OP3 and OP5.  At the author's choice this layer takes the branch Q1b, frame and axis, and asks whether the readout fixes its own
frame: as an extremum of the readout, which is quadratic in \(u\), over all observers, degenerate exactly under boosts along the axis.  Layer r48.2
called the readouts of OP1 and OP5 related, not the same; this layer sharpens that comparison.  Hypothesis H122 was fixed before the computation
(commit \texttt{d861db0e}); H122a and the mode-wise parts of H122b and H122c are not blind, the exact all-order statements and the averages are.
Earlier layers are not rewritten.

\noindent\textbf{Source-only status.}
This v4.10.0-r54 layer is source-only.  It does not build a PDF, does not assign a new DOI,
and does not authorize a public release.  The v5.3.0 Version DOI
\href{https://doi.org/10.5281/zenodo.24434403}{10.5281/zenodo.24434403} and the Concept DOI
\href{https://doi.org/10.5281/zenodo.19210728}{10.5281/zenodo.19210728} are unchanged by it.

\subsection*{v4.10.0-r54 setting}
\label{subsec:v4100-r54-setting}

OP1: a field strength with \(\mathbf E\parallel\mathbf B\) along \(z\) in its rest frame \(R\), read by an observer \(u\) obtained from \(R\) by a boost of rapidity \(\eta\) in an arbitrary
direction at polar angle \(\theta\) to the axis; readout \(Q_F(u)=|\mathbf E'|^2+|\mathbf B'|^2\).  OP5: the transverse tidal readout of r46,
\(\mathcal L=-h\,\Pi^{\mu\nu}u^\alpha u^\beta R^{(1)}_{\mu\alpha\nu\beta}[\gamma]\), with the Fierz--Pauli numerator \(N\) of its source; static modes \(k=(0,k_1,k_2,k_3)\) in the rest
frame \(u_p\) of the preparation; the detector frame \((u,n)\) is the image of \((u_p,e_z)\) under a pure boost of velocity \(b\) in a direction at angle \(t\) to the axis.
Three stages: single modes, the mean over the azimuth of \(k_\perp\), and the isotropic mean over directions of \(k\).

\subsection*{v4.10.0-r54 results}
\label{subsec:v4100-r54-results}

\paragraph{Theoremlet v4.10.0-r54.1 (H122a confirmed: in OP1 the field fixes the frame).}
\(Q_F(u)-Q_F(R)=2(E_0^2+B_0^2)\sin^2\theta\,\sinh^2\eta\ge0\), with equality exactly for boosts along the axis; the Hessian at \(\eta=0\) is \(4(E_0^2+B_0^2)\sin^2\theta\), with kernel exactly
the axis.  The frame \(\mathbf E\parallel\mathbf B\) is the unique minimum of the readout up to boosts along the axis, and the axis is the degenerate direction.
\emph{Proof.} Exact Lorentz transformation of \(\mathbf E,\mathbf B\) (script \texttt{h122}).  \(\square\)

\paragraph{Theoremlet v4.10.0-r54.2 (H122b refuted: in OP5 the preparation frame is a saddle).}
\(N(b,t)-N(0)\) changes sign with the direction of the detector boost in all three stages: for the mode \(k=(1,0,1)\) at \(b=\tfrac12\) it is \(-0.148\) along the
axis and \(+0.296\) across it; in the azimuthal mean (\(a=k_3=1\)) \(-0.148\) and \(+0.139\); in the isotropic mean (\(|k|=1\)) exactly
\(2b^2(5b^2-2)/(45(1-b^2)^2)\) along the axis, negative for \(b^2<\tfrac25\), and \(2b^2(4-b^2)/(45(1-b^2)^2)>0\) across it.
\emph{Proof.} Exact symbolic evaluation with the functions of r46; averages by the angular moments of the quartic numerator.  \(\square\)

\paragraph{Theoremlet v4.10.0-r54.3 (H122c refuted: in OP5 the axis is no degenerate direction).}
Along the axis \(N\) changes for every mode and every azimuthal mean with \(k_3\ne0\) (term proportional to \(b^2k_3^2\)), and in the isotropic mean as in r54.2.
In the isotropic mean no direction keeps \(N=N(0)\) for all \(b\): this would require \(8-24\cos2t=0\) and \(13+24\cos2t+3\cos4t=0\) at once.
\emph{Proof.} Coefficients in \(b\) of the exact isotropic expression \(b^2(8+13b^2-24\cos2t+24b^2\cos2t+3b^2\cos4t)/(180(1-b^2)^2)\).  \(\square\)

\paragraph{Remark v4.10.0-r54.4 (reading; does not carry).}
Against the criterion fixed beforehand H122 does not carry.  In OP1 frame and axis belong to the field: the readout selects them as an extremum.
In OP5 they are an external setting of the detector: the readout does not select them.  This sharpens r48.2 from ``related, not the same'' to a
structural difference.  The assignment named in r48.2 (measuring axis along \(\mathbf E\parallel\mathbf B\), velocity referred to \(R\)) is compatible with the OP5 readout
but not selected by it.  A common law for frame and axis therefore needs an additional setting.  Found after the fact and not part of the
criterion: for a preparation whose modes are all transverse to the axis (\(k_3=0\)), the azimuthal mean gives \(N-N(0)=a^2b^2\sin^2t/3+O(b^4)\), zero exactly along
the axis and positive elsewhere on a grid \(t\in[0,\tfrac\pi2]\), \(b\in[0.1,0.99]\); this preparation has the structure of OP1, with the axis as a property of the
preparation.  It is a new setting that would need its own pre-registered test; the ray along the axis of r32.4 is the opposite case (\(k\parallel n\), \(N=0\)).
The vanishing of the isotropic second-order term at \(\cos^2t=\tfrac23\) is the ordinary quadrupole structure and is not weighted.

\subsection*{v4.10.0-r54 frames}
\label{subsec:v4100-r54-frames}

OP1 in the matter frame M with the field rest frame \(R\); OP5 in the matter frame M with the preparation frame \(u_p\) and the detector frame \((u,n)\).  The two
readouts are compared only in their structure (extremum and degenerate direction), not in their values; no quantity of another frame is used.
Time orientation drops out of both readouts (r46.2, r48.2) and is not tested here (branch Q1a).

\subsection*{v4.10.0-r54 register entries}
\label{subsec:v4100-r54-register}

\begin{center}
\small
\begin{tabular}{@{}>{\raggedright\arraybackslash}p{0.34\textwidth}>{\raggedright\arraybackslash}p{0.58\textwidth}@{}}
\toprule
statement & status after v4.10.0-r54\\
\midrule
H122a (OP1 frame as extremum) & confirmed (not blind)\\
H122b (OP5 preparation frame as extremum) & refuted: saddle in all three stages\\
H122c (OP5 axis as degenerate direction) & refuted\\
common readout law for frame and axis (Q1b) & not supplied by the readout; needs an additional setting\\
OP1, OP5, Q1 & open\\
\bottomrule
\end{tabular}
\end{center}

\subsection*{v4.10.0-r54 notes}
\label{subsec:v4100-r54-notes}

\texttt{llm/\allowbreak Vorab\_\allowbreak Notiz\_\allowbreak H122\_\allowbreak Auslese\_\allowbreak Extremum.md} (commit d861db0e),
\texttt{llm/\allowbreak Ergebnis\_\allowbreak H122\_\allowbreak Auslese\_\allowbreak Extremum.md}, script \texttt{llm/\allowbreak h122\_\allowbreak readout\_\allowbreak extremum.wl},
output \texttt{llm/\allowbreak h122\_\allowbreak readout\_\allowbreak extremum.txt}.

\subsection*{v4.10.0-r54 ledger and guard update}
\label{subsec:v4100-r54-ledger-guard}

\[
R_{\rm chk}^{v4100\text{-}r54}=(1,1,0)
\]
for (status recorded; computations checked in the cited script; a common readout law for frame and axis not derived).

The layer spends no guard:
\[
R_{\rm glob}^{(v4100\text{-}r54,+)}=(0,1,1,1,1),
\qquad
\Gamma_{\rm guard}^{v4100\text{-}r54}=(1,1,1,1,1),
\qquad
\beta_4=1/256\ \text{unspent}.
\]
The layer does not imply exact finite \(C5\) zero closure, exact all-readout terminality, \(C6/T2\), a global photon-clock/RH theorem,
or global all-readout closure.  OP1--OP5, C5 and C6/T2 remain open.
